\documentclass[a4paper, 14pt]{extarticle}

% Поля
%--------------------------------------
\usepackage{geometry}
\geometry{a4paper,tmargin=2cm,bmargin=2cm,lmargin=3cm,rmargin=1cm}
%--------------------------------------


%Russian-specific packages
%--------------------------------------
\usepackage[T2A]{fontenc}
\usepackage[utf8]{inputenc}
\usepackage[english, main=russian]{babel}

%--------------------------------------

\usepackage{textcomp}

% Красная строка
%--------------------------------------
\usepackage{indentfirst}               
%--------------------------------------             


%Graphics
%--------------------------------------
\usepackage{graphicx}
\graphicspath{ {./images/} }
\usepackage{float}%"Плавающие" картинки
%--------------------------------------

% Полуторный интервал
%--------------------------------------
\linespread{1.3}                    
%--------------------------------------

%Выравнивание и переносы
%--------------------------------------
% Избавляемся от переполнений
\sloppy
% Запрещаем разрыв страницы после первой строки абзаца
\clubpenalty=10000
% Запрещаем разрыв страницы после последней строки абзаца
\widowpenalty=10000
%--------------------------------------

%Списки
\usepackage{enumitem}

\usepackage{minted}
%Подписи
\usepackage{caption} 

%Гиперссылки
% \usepackage{hyperref}

\usepackage[colorlinks, urlcolor = blue, filecolor = blue, citecolor = blue, linkcolor = black]{hyperref}


\hypersetup {
unicode=true
}

%Рисунки
%--------------------------------------
\DeclareCaptionLabelSeparator*{emdash}{~--- }
\captionsetup[figure]{labelsep=emdash,font=onehalfspacing,position=bottom}
%--------------------------------------

\usepackage{tempora}

%Листинги
%--------------------------------------
%Листинги
%--------------------------------------
\usepackage{listings}
\lstset{
  basicstyle=\ttfamily\footnotesize, 
  %basicstyle=\footnotesize\AnkaCoder,        % the size of the fonts that are used for the code
  breakatwhitespace=false,         % sets if automatic breaks shoulbd only happen at whitespace
  breaklines=true,                 % sets automatic line breaking
  captionpos=t,                    % sets the caption-position to bottom
  inputencoding=utf8,
  frame=single,                    % adds a frame around the code
  keepspaces=true,                 % keeps spaces in text, useful for keeping indentation of code (possibly needs columns=flexible)
  keywordstyle=\bf,       % keyword style
  numbers=left,                    % where to put the line-numbers; possible values are (none, left, right)
  numbersep=5pt,                   % how far the line-numbers are from the code
  xleftmargin=25pt,
  xrightmargin=25pt,
  showspaces=false,                % show spaces everywhere adding particular underscores; it overrides 'showstringspaces'
  showstringspaces=false,          % underline spaces within strings only
  showtabs=false,                  % show tabs within strings adding particular underscores
  stepnumber=1,                    % the step between two line-numbers. If it's 1, each line will be numbered
  tabsize=2,                       % sets default tabsize to 8 spaces
  title=\lstname                   % show the filename of files included with \lstinputlisting; also try caption instead of title
}
%--------------------------------------

%%% Математические пакеты %%%
%--------------------------------------
\usepackage{amsthm,amsfonts,amsmath,amssymb,amscd}  % Математические дополнения от AMS
\usepackage{mathtools}                              % Добавляет окружение multlined
\usepackage[perpage]{footmisc}


% графики
\usepackage{booktabs}        % Пакет для улучшенных таблиц
\usepackage{pgfplots}        % Пакет для построения графиков
\pgfplotsset{compat=1.17}    % Установка совместимости

% Литература 

% переименовываем  список литературы в "список используемой литературы"
\addto\captionsrussian{\def\refname{Список литературы}}



\captionsetup[figure]{name=Рисунок, labelsep=space}


%--------------------------------------

%--------------------------------------
%			НАЧАЛО ДОКУМЕНТА
%--------------------------------------

\begin{document}

%--------------------------------------
%			ТИТУЛЬНЫЙ ЛИСТ
%--------------------------------------
\begin{titlepage}
\thispagestyle{empty}
\newpage


%Шапка титульного листа
%--------------------------------------
\vspace*{-60pt}
\hspace{-65pt}
\begin{minipage}{0.3\textwidth}
\hspace*{-20pt}\centering
\includegraphics[width=\textwidth]{emblem}
\end{minipage}
\begin{minipage}{0.67\textwidth}\small \textbf{
\vspace*{-0.7ex}
\hspace*{-6pt}\centerline{Министерство науки и высшего образования Российской Федерации}
\vspace*{-0.7ex}
\centerline{Федеральное государственное автономное образовательное учреждение }
\vspace*{-0.7ex}
\centerline{высшего образования}
\vspace*{-0.7ex}
\centerline{<<Московский государственный технический университет}
\vspace*{-0.7ex}
\centerline{имени Н.Э. Баумана}
\vspace*{-0.7ex}
\centerline{(национальный исследовательский университет)>>}
\vspace*{-0.7ex}
\centerline{(МГТУ им. Н.Э. Баумана)}}
\end{minipage}
%--------------------------------------

%Полосы
%--------------------------------------
\vspace{-25pt}
\hspace{-35pt}\rule{\textwidth}{2.3pt}

\vspace*{-20.3pt}
\hspace{-35pt}\rule{\textwidth}{0.4pt}
%--------------------------------------

\vspace{1.5ex}
\hspace{-35pt} \noindent \small ФАКУЛЬТЕТ\hspace{80pt} <<Информатика и системы управления>>

\vspace*{-16pt}
\hspace{47pt}\rule{0.83\textwidth}{0.4pt}

\vspace{0.5ex}
\hspace{-35pt} \noindent \small КАФЕДРА\hspace{50pt} <<Теоретическая информатика и компьютерные технологии>>

\vspace*{-16pt}
\hspace{30pt}\rule{0.866\textwidth}{0.4pt}

\vspace{11em}

\begin{center}
\Large {\bf Домашнее задание} \\ 
\large <<Решение баллистической задачи. Модель Ньютона.>> \\
\large {по курсу <<Моделирование>>} \\

\end{center}\normalsize

\vspace{8em}

\vbox{%
\hfill%
\vbox{%
\hbox{Студент: Аринин М.А.}%
\hbox{Группа: ИУ9-82Б}%
\hbox{Преподаватель: Домрачева А.Б.}%
}%
} 

\bigskip

\vfill


\begin{center}
\textsl{Москва 2026}
\end{center}
\end{titlepage}
%--------------------------------------
%		КОНЕЦ ТИТУЛЬНОГО ЛИСТА
%--------------------------------------


\renewcommand{\ttdefault}{pcr}

\setlength{\tabcolsep}{3pt}

\newpage
\setcounter{page}{2}

\section{Постановка задачи}
Рассматривается задача о движении тела, брошенного под углом к горизонту. Тело массой $m$ запускается с начальной скоростью $v_0$ под углом $\alpha$ к горизонту и движется в вертикальной плоскости под действием силы тяжести. В зависимости от учета сопротивления воздуха различают две модели: модель Галилея (без учета сопротивления) и модель Ньютона (с учетом сопротивления воздуха, пропорционального квадрату скорости).

Цель работы --- численное решение данной задачи методом Рунге-Кутты 4-го порядка для обеих моделей, сравнение траекторий и анализ влияния сопротивления воздуха на движение тела.

\textbf{Входные параметры:}
\begin{itemize}
  \item $g = 9.81$ м/с² — ускорение свободного падения
  \item $v_0 = 60.0$ м/с — начальная скорость
  \item $\alpha = 43°$ — угол бросания
  \item $C = 0.15$ — коэффициент сопротивления
  \item $\rho_{\text{air}} = 1.225$ кг/м³ — плотность воздуха
  \item $r_{\text{core}} = 0.04$ м — радиус ядра
  \item $\rho_{\text{core}} = 7300.0$ кг/м³ — плотность олова
\end{itemize}

\textbf{Выходные параметры:}
\begin{itemize}
  \item $S = \pi r_{\text{core}}^2$ — площадь поперечного сечения
  \item $V_{\text{core}} = \frac{4}{3}\pi r_{\text{core}}^3$ — объем ядра
  \item $m = \rho_{\text{core}} \cdot V_{\text{core}}$ — масса ядра
  \item $\beta = \frac{C \rho_{\text{air}} S}{2}$ — коэффициент сопротивления
  \item $R$ — дальность полета
\end{itemize}

%--------------------------------------------------------------------
\section{Построение математической модели}
Рассматривается движение тела в вертикальной плоскости $(x,y)$ при следующих допущениях:
\begin{itemize}
  \item ускорение свободного падения $g$ постоянно;
  \item плотность воздуха $\rho_{\text{air}}$ постоянна;
  \item тело рассматривается как материальная точка массой $m$;
  \item сопротивление воздуха направлено против скорости и пропорционально квадрату модуля скорости.
\end{itemize}

Введем переменные состояния: координаты $(x,y)$ и проекции скорости $(u,w)$.

\subsection{Модель Галилея}

\begin{figure}[H]
    \centering
    \includegraphics[width=1\textwidth]{1.jpg}
    \label{fig:galilei}
  \end{figure}

В модели Галилея сопротивление воздуха не учитывается. Система дифференциальных уравнений имеет вид:
\[
\begin{cases}
  \dfrac{dx}{dt} = u,\\
  \dfrac{dy}{dt} = w,\\
  \dfrac{du}{dt} = 0,\\
  \dfrac{dw}{dt} = -g.
\end{cases}
\]

Начальные условия:
\[
x(0)=0,\quad y(0)=0,\quad u(0)=v_0\cos\alpha,\quad w(0)=v_0\sin\alpha.
\]

В данной модели траектория движения является параболой. Аналитическое решение:
\[
x(t) = v_0\cos\alpha \cdot t, \quad y(t) = v_0\sin\alpha \cdot t - \frac{g}{2}t^2.
\]

\subsection{Модель Ньютона}
\begin{figure}[H]
    \centering
    \includegraphics[width=1\textwidth]{2.jpg}
    \label{fig:newton}
  \end{figure}
В модели Ньютона учитывается сопротивление воздуха, пропорциональное квадрату скорости. Сила сопротивления задается как
\[
\vec{F}_d = -\beta V \vec{v}, \qquad V = \sqrt{u^2+w^2},
\]
где $\beta = \dfrac{C\,\rho_{\text{air}}\,S}{2}$.

Система дифференциальных уравнений:
\[
\begin{cases}
  \dfrac{dx}{dt} = u,\\
  \dfrac{dy}{dt} = w,\\
  \dfrac{du}{dt} = -\dfrac{\beta}{m}uV,\\
  \dfrac{dw}{dt} = -g-\dfrac{\beta}{m}wV.
\end{cases}
\]

Начальные условия:
\[
x(0)=0,\quad y(0)=0,\quad u(0)=v_0\cos\alpha,\quad w(0)=v_0\sin\alpha.
\]

Связь параметров с геометрией и материалом для сферического ядра радиуса $r_{\text{core}}$:
\[
S=\pi r_{\text{core}}^2,\qquad V_{\text{core}}=\frac{4}{3}\pi r_{\text{core}}^3,\qquad m=\rho_{\text{core}}V_{\text{core}},
\qquad \beta=\frac{C\rho_{\text{air}}S}{2}.
\]

В данной модели траектория движения не является параболой из-за нелинейного члена сопротивления.

\section{Решение вычислительной задачи}
Численное решение строится в виде последовательности подзадач.

\subsection{Постановка задачи Коши}
Система ОДУ первого порядка с начальными условиями представляет собой корректную задачу Коши. Для модели Ньютона правая часть непрерывна. При $V\to 0$ возможна численная неустойчивость (деление на очень малые значения), поэтому в реализации вводится проверка $V<\varepsilon$.

\subsection{Выбор временного интервала}
Для модели Галилея высота
\[
y(t)=v_0\sin\alpha\cdot t-\frac{g}{2}t^2
\]
обращается в ноль при двух значениях времени: $t_1=0$ (вырожденное решение в начальной точке) и
\[
t_2=\frac{2v_0\sin\alpha}{g}
\]
(время полета). Для численного расчета берется невырожденный корень $t_2$ и небольшой запас по времени.

\subsection{Численное интегрирование}
Используется метод Рунге--Кутты 4-го порядка (RK4):
\[
\begin{aligned}
k_1 &= \Delta t\, f(t_n, y_n),\\
k_2 &= \Delta t\, f\!\left(t_n+\frac{\Delta t}{2},\, y_n+\frac{k_1}{2}\right),\\
k_3 &= \Delta t\, f\!\left(t_n+\frac{\Delta t}{2},\, y_n+\frac{k_2}{2}\right),\\
k_4 &= \Delta t\, f(t_n+\Delta t,\, y_n+k_3),\\
y_{n+1} &= y_n+\frac{k_1+2k_2+2k_3+k_4}{6}.
\end{aligned}
\]

\subsection{Критерий окончания полета}
Интегрирование останавливается при $y<0$ (касание земли). Дальность полета оценивается как $R=\max x(t)$ на интервале, где $y(t)\ge 0$.

\section{Анализ моделей}
Сравнение двух моделей позволяет оценить влияние сопротивления воздуха на траекторию движения тела.

В модели Галилея траектория параболическая и дальность максимальна при заданных $v_0$ и $\alpha$. В модели Ньютона дальность уменьшается, а траектория становится непараболической из-за нелинейного члена сопротивления. Влияние сопротивления воздуха усиливается при росте $v_0$, увеличении $r_{\text{core}}$ (рост $S$) и уменьшении плотности материала (уменьшение $m$ при фиксированном $r_{\text{core}}$).

Ограничения моделей: постоянное $g$, отсутствие ветра, постоянные $C$ и $\rho_{\text{air}}$, отсутствие подъемной силы и вращения тела. Эти допущения упрощают расчет, но ограничивают применимость результатов.

\begin{figure}[H]
  \centering
  \includegraphics[width=0.9\textwidth]{trajectory_galilei.png}
  \caption{Траектория движения тела в модели Галилея}
  \label{fig:analysis_galilei}
\end{figure}

\begin{figure}[H]
  \centering
  \includegraphics[width=0.9\textwidth]{trajectory_newton.png}
  \caption{Траектория движения тела в модели Ньютона}
  \label{fig:analysis_newton}
\end{figure}

\begin{figure}[H]
  \centering
  \includegraphics[width=0.9\textwidth]{trajectory_comparison.png}
  \caption{Сравнение траекторий моделей Галилея и Ньютона}
  \label{fig:analysis_comparison}
\end{figure}
\section{Практическая реализация}
\newenvironment{longlisting}{\captionsetup{type=listing}}{}
\begin{longlisting}
\caption{Модель Галилея и Ньютона}
\begin{minted}[frame=single,fontsize = \footnotesize, linenos, xleftmargin = 1.5em, breaklines=true]{python}
import numpy as np
import matplotlib.pyplot as plt


class RungeKuttaSolver:
    def __init__(self, dt=0.01):
        self.dt = dt
    
    def rk4_step(self, f, t, y):
        k1 = self.dt * f(t, y)
        k2 = self.dt * f(t + self.dt/2, y + k1/2)
        k3 = self.dt * f(t + self.dt/2, y + k2/2)
        k4 = self.dt * f(t + self.dt, y + k3)
        
        return y + (k1 + 2*k2 + 2*k3 + k4) / 6
    
    def solve(self, f, y0, t_span):
        y = np.zeros((len(t_span), len(y0)))
        y[0] = y0
        
        for i in range(1, len(t_span)):
            y[i] = self.rk4_step(f, t_span[i-1], y[i-1])
            
        return y


class GalileiModel:
    def __init__(self, g=9.81, v0=1.0, alpha=45.0):
        self.g = g
        self.v0 = v0
        self.alpha = np.deg2rad(alpha)
    
    def solve_numerically(self, t_span, dt=0.01):
        t_flight = 2.0 * self.v0 * np.sin(self.alpha) / self.g
        

        if len(t_span) > 0 and t_span[-1] < t_flight * 1.5:
            t_max = t_flight * 1.5
            t_span = np.arange(0, t_max, dt)
        x0 = 0.0
        y0 = 0.0
        u0 = self.v0 * np.cos(self.alpha)
        w0 = self.v0 * np.sin(self.alpha)
        
        y_init = np.array([x0, y0, u0, w0])
        
        def f_galilei(t, y):
            x, y_pos, u, w = y
            return np.array([u, w, 0, -self.g])
        
        solver = RungeKuttaSolver(dt=dt)
        solution = solver.solve(f_galilei, y_init, t_span)
        
        x = solution[:, 0]
        y = solution[:, 1]

        mask = y >= 0
        if not np.all(mask):
            first_negative = np.where(~mask)[0]
            if len(first_negative) > 0:
                mask = np.arange(len(y)) < first_negative[0]
        else:
            last_positive = np.where(y >= 0)[0]
            if len(last_positive) > 0:
                mask = np.arange(len(y)) <= last_positive[-1]
        
        return x[mask], y[mask]


class NewtonModel:
    def __init__(self, g=9.81, v0=1.0, alpha=45.0, m=None, r_core=None, rho_core=None, 
                 beta=None, C=0.15, rho_air=1.225, S=None):
        self.g = g
        self.v0 = v0
        self.alpha = np.deg2rad(alpha)
        
        if m is None:
            V_core = (4.0 / 3.0) * np.pi * r_core**3
            self.m = rho_core * V_core
            self.r_core = r_core
            self.rho_core = rho_core
        else:
            self.m = m
            self.r_core = r_core
            self.rho_core = rho_core
        
        if S is None:
            if r_core is not None:
                self.S = np.pi * r_core**2
        else:
            self.S = S
        
        if beta is None:
            self.beta = C * rho_air * self.S / 2.0
        else:
            self.beta = beta
    
    def solve(self, t_span, dt=0.01):
        x0 = 0.0
        y0 = 0.0
        u0 = self.v0 * np.cos(self.alpha)
        w0 = self.v0 * np.sin(self.alpha)
        
        y_init = np.array([x0, y0, u0, w0])
        
        def f_newton(t, y):
            x, y_pos, u, w = y
            V = np.sqrt(u**2 + w**2)

            if V < 1e-10:
                du_dt = 0
                dw_dt = -self.g
            else:
                du_dt = -self.beta * u * V / self.m
                dw_dt = -self.g - self.beta * w * V / self.m
            
            return np.array([u, w, du_dt, dw_dt])
        
        solver = RungeKuttaSolver(dt=dt)
        solution = solver.solve(f_newton, y_init, t_span)
        
        x = solution[:, 0]
        y = solution[:, 1]
        u = solution[:, 2]
        w = solution[:, 3]
        
        mask = y >= 0
        return x[mask], y[mask], u[mask], w[mask]


def plot_trajectories(g=9.81, v0=30.0, alpha=45.0, C=0.5, rho_air=1.225, 
                      S=None, r_core=0.04, rho_core=200.0, diff_range=None, dt=0.001, 
                      filename='trajectory_comparison.png'):
    if S is None:
        S = np.pi * r_core**2
    
    t_flight_galilei = 2.0 * v0 * np.sin(np.deg2rad(alpha)) / g
    t_max = max(5.0, t_flight_galilei * 1.5)
    t_span = np.arange(0, t_max, dt)
    
    galilei = GalileiModel(g=g, v0=v0, alpha=alpha)
    x_gal, y_gal = galilei.solve_numerically(t_span, dt=dt)
    
    newton = NewtonModel(g=g, v0=v0, alpha=alpha, r_core=r_core, rho_core=rho_core, 
                         C=C, rho_air=rho_air)
    x_newt, y_newt, u_newt, w_newt = newton.solve(t_span, dt=dt)
    
    plt.figure(figsize=(8, 6))
    if len(x_gal) > 0 and len(y_gal) > 0:
        plt.plot(x_gal, y_gal, color='#0066CC', linestyle='-', linewidth=2.5, 
                label='Галилей', zorder=2)
        x_max_gal = np.max(x_gal)
        y_max_gal = np.max(y_gal)
        plt.xlim(-0.005, x_max_gal * 1.05)
        plt.ylim(-0.005, y_max_gal * 1.1)
        
        plt.axhline(y=0, color='gray', linestyle='--', linewidth=1, alpha=0.5, zorder=0)
        plt.axvline(x=0, color='gray', linestyle='--', linewidth=1, alpha=0.5, zorder=0)
    
    plt.xlabel('x (м)', fontsize=11)
    plt.ylabel('y (м)', fontsize=11)
    plt.title('Траектория Галилея\n(упрощенный Ньютон без сопротивления воздуха)', 
              fontsize=12, fontweight='bold')
    plt.legend(fontsize=9, loc='upper right', framealpha=0.9)
    plt.grid(True, alpha=0.3, linestyle=':')
    plt.axis('equal')
    plt.tight_layout()
    plt.savefig('trajectory_galilei.png', dpi=300, bbox_inches='tight')
    plt.close()
    
    plt.figure(figsize=(8, 6))
    if len(x_newt) > 0 and len(y_newt) > 0:
        x_parabola = np.linspace(0, np.max(x_newt), 1000)
        y_parabola = -g / (2 * v0**2 * np.cos(np.deg2rad(alpha))**2) * x_parabola**2 + np.tan(np.deg2rad(alpha)) * x_parabola
        y_parabola = y_parabola[y_parabola >= 0]
        x_parabola = x_parabola[:len(y_parabola)]
        
        plt.plot(x_parabola, y_parabola, color='#00CCFF', linestyle='--', linewidth=1.5, 
                label='Парабола (для сравнения)', zorder=1, alpha=0.6)
        plt.plot(x_newt, y_newt, color='#CC0000', linestyle='-', linewidth=2.5, 
                label='Ньютон (непараболическая)', zorder=2)
        x_max_newt = np.max(x_newt)
        y_max_newt = np.max(y_newt)
        plt.xlim(-0.005, x_max_newt * 1.05)
        plt.ylim(-0.005, y_max_newt * 1.1)
        
        plt.axhline(y=0, color='gray', linestyle='--', linewidth=1, alpha=0.5, zorder=0)
        plt.axvline(x=0, color='gray', linestyle='--', linewidth=1, alpha=0.5, zorder=0)
    
    plt.xlabel('x (м)', fontsize=11)
    plt.ylabel('y (м)', fontsize=11)
    plt.title('Траектория Ньютона\n(с учетом сопротивления воздуха)', 
              fontsize=12, fontweight='bold')
    plt.legend(fontsize=9, loc='upper right', framealpha=0.9)
    plt.grid(True, alpha=0.3, linestyle=':')
    plt.axis('equal')
    plt.tight_layout()
    plt.savefig('trajectory_newton.png', dpi=300, bbox_inches='tight')
    plt.close()
    
    plt.figure(figsize=(8, 6))
    if len(x_gal) > 0 and len(x_newt) > 0:
        if diff_range is None:
            range_gal = np.max(x_gal)
            range_newt = np.max(x_newt)
            diff_range = range_gal - range_newt
        else:
            range_gal = np.max(x_gal)
            range_newt = np.max(x_newt)
        
        t_gal = np.arange(0, len(x_gal)) * dt
        t_newt = np.arange(0, len(x_newt)) * dt
        
        t_min = 0
        t_max_common = min(t_gal[-1] if len(t_gal) > 0 else 0, t_newt[-1] if len(t_newt) > 0 else 0)
        t_common = np.linspace(t_min, t_max_common, 1000)
        
        x_gal_interp = np.interp(t_common, t_gal, x_gal)
        x_newt_interp = np.interp(t_common, t_newt, x_newt)
        
        diff_x = x_gal_interp - x_newt_interp
        
        plt.plot(x_gal_interp, diff_x, color='purple', linestyle='-', linewidth=2.5, 
                label='Разница по x (Галилей - Ньютон)', zorder=2)
        plt.axhline(y=0, color='black', linestyle='--', linewidth=1, alpha=0.5)
        
        plt.xlabel('x (м)', fontsize=11)
        plt.ylabel('Разница x (м)', fontsize=11)
        plt.title('Погрешность по дальности (x)', fontsize=12, fontweight='bold')
        plt.legend(fontsize=9, loc='best', framealpha=0.9)
        plt.grid(True, alpha=0.3, linestyle=':')
    
    plt.tight_layout()
    plt.savefig('trajectory_error.png', dpi=300, bbox_inches='tight')
    plt.close()
    
    plt.figure(figsize=(8, 6))
    if len(x_gal) > 0 and len(x_newt) > 0:
        plt.plot(x_gal, y_gal, color='#0066CC', linestyle='-', linewidth=2.5, 
                label='Галилей', zorder=2)
        plt.plot(x_newt, y_newt, color='#CC0000', linestyle='-', linewidth=2.5, 
                label='Ньютон', zorder=3)
        plt.axhline(y=0, color='gray', linestyle='--', linewidth=1, alpha=0.5, zorder=0)
        plt.axvline(x=0, color='gray', linestyle='--', linewidth=1, alpha=0.5, zorder=0)
        
        x_max = max(np.max(x_gal), np.max(x_newt))
        y_max = max(np.max(y_gal), np.max(y_newt))
        plt.xlim(-0.005, x_max * 1.05)
        plt.ylim(-0.005, y_max * 1.1)
    
    plt.xlabel('x (м)', fontsize=11)
    plt.ylabel('y (м)', fontsize=11)
    plt.title('Сравнение траекторий', fontsize=12, fontweight='bold')
    plt.legend(fontsize=9, loc='upper right', framealpha=0.9)
    plt.grid(True, alpha=0.3, linestyle=':')
    
    plt.tight_layout()
    plt.savefig('trajectory_comparison.png', dpi=300, bbox_inches='tight')
    plt.close()


if __name__ == "__main__":
    g = 9.81
    v0 = 60.0
    alpha = 43.0
    C = 0.15
    rho_air = 1.225
    r_core = 0.04
    rho_core = 7300.0
    
    S = np.pi * r_core**2
    V_core = (4.0 / 3.0) * np.pi * r_core**3
    m = rho_core * V_core
    beta = C * rho_air * S / 2.0
    
    t_flight_galilei = 2.0 * v0 * np.sin(np.deg2rad(alpha)) / g
    t_max = max(1.0, t_flight_galilei * 1.5)
    dt = 0.001
    t_span = np.arange(0, t_max, dt)
    
    print("=" * 60)
    print("ВХОДНЫЕ ПАРАМЕТРЫ")
    print("=" * 60)
    print(f"  g = {g} m/s^2")
    print(f"  v0 = {v0} m/s")
    print(f"  alpha = {alpha} grad")
    print(f"  C = {C}")
    print(f"  rho_air = {rho_air} kg/m^3")
    print(f"  S = {S} m^2 ({S*100:.1f} dm^2 = {S*10000:.0f} cm^2)")
    print(f"  rho_core = {rho_core} kg/m^3")
    
    print("\n" + "=" * 60)
    print("ВЫЧИСЛЯЕМЫЕ ПАРАМЕТРЫ")
    print("=" * 60)
    print(f"  r_core = sqrt(S/pi) = {r_core:.6f} m")
    print(f"  V_core = (4/3)*pi*r^3 = {V_core:.6f} m^3")
    print(f"  m = rho_core * V_core = {m:.6f} kg")
    print(f"  beta = C*rho_air*S/2 = {beta:.6f}")
    print(f"  t_flight (Galilei) = 2*v0*sin(alpha)/g = {t_flight_galilei:.4f} s")
    print(f"  t_max = {t_max:.4f} s")
    print(f"  dt = {dt} s")
    
    print("\n" + "=" * 60)
    print("Модель Галилея (упрощенный Ньютон без сопротивления воздуха)")
    print("=" * 60)
    galilei = GalileiModel(g=g, v0=v0, alpha=alpha)
    x_gal, y_gal = galilei.solve_numerically(t_span, dt=dt)
    range_gal = np.max(x_gal)
    print(f"Дальность полёта: {range_gal:.4f} m")
    
    print("\n" + "=" * 60)
    print("Модель Ньютона (с учетом сопротивления воздуха)")
    print("=" * 60)
    newton = NewtonModel(g=g, v0=v0, alpha=alpha, r_core=r_core, rho_core=rho_core, 
                         C=C, rho_air=rho_air)
    x_newt, y_newt, u_newt, w_newt = newton.solve(t_span, dt=dt)
    range_newt = np.max(x_newt)
    print(f"Дальность полёта: {range_newt:.4f} m")

    print("\n" + "=" * 60)
    print("СРАВНЕНИЕ ДАЛЬНОСТИ ПОЛЁТА")
    print("=" * 60)
    diff_range = range_gal - range_newt
    
    print(f"Разница в дальности: {diff_range:.6f} m ({diff_range*1000:.3f} mm)")
    
    relative_error = abs(diff_range / range_gal) * 100
    print(f"Относительная погрешность: {relative_error:.6f} %")
    

    plot_trajectories(g=g, v0=v0, alpha=alpha, C=C, rho_air=rho_air, 
                     S=S, r_core=r_core, rho_core=rho_core, 
                     diff_range=diff_range, dt=dt)

\end{minted}
\end{longlisting}

\section{Результаты}

\begin{figure}[H]
  \centering
    \includegraphics[width=1\textwidth]{image.png}
    \label{fig:newton}
\end{figure}

В результате выполнения программы получены следующие результаты:

\begin{itemize}
  \item \textbf{Дальность полета методом Галилея:} вычисляется численно методом Рунге-Кутты 4-го порядка
  \item \textbf{Дальность полета методом Ньютона:} вычисляется численно с учетом сопротивления воздуха
  \item \textbf{Разница в дальности:} показывает влияние сопротивления воздуха на траекторию
  \item \textbf{Графики траекторий:} построены отдельные графики для каждой модели и график сравнения
\end{itemize}

%--------------------------------------------------------------------
\section{Выводы}
\begin{itemize}
  \item Метод Галилея дает параболическую траекторию, что соответствует классической теории движения тела под действием силы тяжести без учета сопротивления воздуха.
  
  \item Метод Ньютона с учетом сопротивления воздуха показывает, что траектория \emph{не является параболой}. Сопротивление воздуха, пропорциональное квадрату скорости ($F = -\beta V^2$), приводит к уменьшению дальности полета и изменению формы траектории.
  
  \item В упрощенной модели (плоская Земля, постоянное $g = 9.81$ м/с²) траектория Ньютона не является ни параболой (из-за сопротивления), ни эллипсом/гиперболой (для этого требуется центральное поле тяготения с переменным $g$). Это более сложная непараболическая кривая.
  
  \item Разница в дальности полета между методами зависит от параметров задачи: начальной скорости, массы тела, площади поперечного сечения и коэффициента сопротивления. При малой массе и большой площади сечения влияние сопротивления воздуха становится значительным.
  
  \item Метод Рунге-Кутты 4-го порядка обеспечивает высокую точность численного интегрирования системы дифференциальных уравнений и позволяет получить траектории с достаточной точностью для анализа различий между методами.
\end{itemize}

\end{document}